\chapter[Аштавакра]{Настанови щодо самопізнання (Аштавакра)}

[1.2] Щоб бути цілком вільним, уникай полону чуттєвих насолод так, як уникаєш отрути. Плекай терпимість, щирість, співчуття, простоту й чесність.

[1.3] Ти --- не земля, не вода, не вогонь і не повітря. І навіть не порожній простір. Звільнення --- це усвідомлення себе чистою Свідомістю, свідком усього цього.

[1.4] Якщо зможеш побачити, що ти --- чиста Свідомість, і перестанеш ототожнювати себе з тілом, миттєво здобудеш щастя, спокій і звільнення.

[1.5] Ти не належиш до жодної касти й не маєш визначених обов'язків. Твоя природа невидима для очей, вільна від прив'язаностей і безформна. Вона --- свідок усього. Будь щасливий.

[1.6] Правильне й неправильне, насолода й біль перебувають лише в розумі та не стосуються тебе. Не ти дієш і не ти зазнаєш наслідків, тож ти завжди вільний.

[1.7] Ти --- єдиний Свідок, і ти завжди вільний. Причина твого поневолення --- у тому, що ти вважаєш себе чимось іншим.

[1.8] Думка «я дію» --- укус отруйної змії. Знання «я нічого не роблю» --- мудрість-протиотрута. Будь щасливий.

[1.9] Саме усвідомлення «я --- єдина Свідомість» спалює все невігластво й страждання у вогні однієї миті. Будь щасливий.

[1.10] Ти --- безмежна Свідомість, блаженство, найвищий екстаз, у якій Всесвіт набуває форми, наче дивовижна змія, що ввижається в мотузці. Будь щасливий.

[1.11] Правду кажуть: «думка формує твою реальність». Якщо вважаєш себе обмеженим --- ти обмежений. Якщо вважаєш себе безмежним --- ти безмежний.

[1.12] Твоя істинна природа досконала, вільна, всюдисуща, єдина, ні з чим не ототожнена, позбавлена бажань і сповнена спокою. Сансара --- ілюзія, породжена невіглаством. Це існування без самопізнання.

[1.13] Споглядай себе лише як Свідомість --- саму Єдність. Відмовся від думки, ніби ти відокремлений, ніби свідомість блукає сама. Відмовся від поділу на внутрішнє й зовнішнє.

[1.14] Уже давно ти перебуваєш в пастці ототожнення з тілом, думаючи «я --- особа». Нехай знання «я --- лише Свідомість» стане мечем, що визволить тебе. Будь щасливий.

[1.15] Ти завжди був вільним, осяйним і незаплямованим. Намагання силоміць заспокоїти розум --- саме те, що тримає тебе в неволі.

[1.16] Ти --- Свідомість, сутність Всесвіту. Реальність існує завдяки тобі. Не обмежуй свого мислення.

[1.17] Ти нічим не зумовлений, незмінний, безформний, нерухомий, незбагненний і спокійний. Нічого не бажай. Просто будь чистою Свідомістю.

[1.18] Форма здається реальною, тоді як безформне --- вічне. Усвідомивши це, ти більше не впадеш в оману.

[1.19] Як у дзеркалі образ і відбитий предмет існують усередині й зовні, так і Вище Я існує всередині тіла й поза ним.

[1.20] Як всюдисущий простір існує і всередині, і зовні глечика, так позачасова Єдність пронизує все.
